% ============================================================
%  Return-Conditioned Emergency Corridor Optimization
%  via Decision Transformer
%
%  Target venue: IEEE Transactions on Intelligent
%  Transportation Systems (T-ITS)
% ============================================================
\documentclass[journal]{IEEEtran}

% ---- Packages ---------------------------------------------
\usepackage[utf8]{inputenc}
\usepackage[T1]{fontenc}
\usepackage{amsmath,amssymb,amsfonts}
\usepackage{algorithmic}
\usepackage{algorithm}
\usepackage{graphicx}
\usepackage{textcomp}
\usepackage{xcolor}
\usepackage{booktabs}
\usepackage{multirow}
\usepackage{hyperref}
\usepackage{cleveref}
\usepackage{subcaption}
\usepackage{bm}

% ---- TODO command -----------------------------------------
\newcommand{\todo}[1]{\textcolor{red}{\textbf{[TODO: #1]}}}

% ---- Math shortcuts ---------------------------------------
\newcommand{\R}{\mathbb{R}}
\newcommand{\E}{\mathbb{E}}
\newcommand{\Rtg}{\widehat{R}}
\DeclareMathOperator*{\argmax}{arg\,max}
\DeclareMathOperator*{\argmin}{arg\,min}

% ============================================================
\begin{document}

\title{Return-Conditioned Emergency Corridor Optimization\\
via Decision Transformer}

\author{Haoran~Su%
\thanks{Manuscript received XXXX XX, 2026; revised XXXX XX, 2026.}%
\thanks{H.~Su is with \todo{affiliation, department, university,
city, country} (e-mail: \todo{email}).}}

\markboth{IEEE Transactions on Intelligent Transportation
Systems}{Su: Return-Conditioned Emergency Corridor Optimization
via Decision Transformer}

\maketitle

% ============================================================
% ABSTRACT
% ============================================================
\begin{abstract}
Emergency vehicle (EV) response time is a critical determinant of
life-saving outcomes, yet current signal preemption strategies remain
largely reactive, triggering phase changes only when an EV is detected
within a short upstream range.  We propose a \emph{return-conditioned}
framework for joint route and signal corridor optimization based on the
Decision Transformer (DT) architecture.  By casting the corridor
optimization problem as offline, return-conditioned sequence modeling,
our approach (1)~eliminates the need for online environment interaction
during training, (2)~enables dispatch-level control of urgency through
a single scalar target-return knob, and (3)~extends naturally to
multi-agent settings via a Multi-Agent Decision Transformer (MADT) with
graph attention layers for spatial coordination across intersections.
Comprehensive experiments on the LightSim simulator over grid and
real-world-inspired networks demonstrate that the proposed method
reduces average EV travel time by \todo{XX\%} relative to fixed-timing
preemption and matches or exceeds online RL baselines (PPO, DQN) while
requiring no exploratory interaction with the traffic network.

\todo{Refine numbers once experiments are complete.}
\end{abstract}

\begin{IEEEkeywords}
Decision Transformer, emergency vehicle preemption, offline
reinforcement learning, multi-agent reinforcement learning, traffic
signal control, return conditioning.
\end{IEEEkeywords}

% ============================================================
% I. INTRODUCTION
% ============================================================
\section{Introduction}
\label{sec:introduction}

Every second of delay in emergency vehicle (EV) response correlates
with measurably worse patient outcomes in medical emergencies and
greater property loss in fire incidents~\cite{qin2012control}.
Modern urban traffic signal systems offer \emph{signal preemption}
---the ability to override normal phase plans to create a green
corridor for an approaching EV---but the vast majority of deployed
preemption strategies are \emph{reactive}: they detect the EV via
optical or radio sensors only a few hundred meters upstream and force
an immediate phase change, often without regard for downstream
congestion, queue spillback, or alternative routing
\cite{wang2022ev,haydari2020deep}.

Recent advances in deep reinforcement learning (RL) have shown promise
for adaptive traffic signal control~\cite{wei2019presslight,
wei2019colight,chen2020toward,chu2019multi}, and several works have
applied online RL to EV preemption~\cite{wang2022ev}.  However, online
RL methods face two practical barriers in safety-critical EV settings:
\begin{enumerate}
  \item \textbf{Exploratory risk.}  Online algorithms such as
    PPO~\cite{schulman2017proximal} and DQN~\cite{mnih2015human}
    require millions of environment steps during which suboptimal
    exploration may endanger real EVs or surrounding traffic.
  \item \textbf{No intuitive control interface.}  Once trained, an RL
    policy offers a fixed mapping from state to action with no
    straightforward mechanism for a dispatcher to request ``slightly
    faster but more disruptive'' versus ``moderate speed with minimal
    civilian delay.''
\end{enumerate}

The Decision Transformer (DT)~\cite{chen2021decision} addresses both
issues by reformulating RL as \emph{return-conditioned sequence
modeling}.  A DT is trained purely offline on previously collected
trajectories and, at inference time, conditions on a
\emph{target return} to modulate policy behavior without retraining.
This property is especially attractive for EV corridor optimization,
where dispatch operators could dial a target return that trades off EV
speed against civilian disruption in real time.

Despite these advantages, no prior work has applied the Decision
Transformer paradigm---or, more broadly, offline sequence-modeling
RL---to the problem of joint route and signal corridor optimization
for emergency vehicles.

\subsection*{Contributions}

This paper makes the following contributions:
\begin{enumerate}
  \item \textbf{First application of Decision Transformer to emergency
    vehicle corridor optimization.}  We formalize the EV corridor
    problem as a return-conditioned sequence modeling task and show
    that a DT trained entirely offline on mixed-quality trajectory data
    can learn effective signal-preemption and routing strategies
    without any online interaction with the traffic network.

  \item \textbf{Return-conditioning as a dispatch-level control
    interface for adjustable urgency.}  We demonstrate that varying
    the target return at inference time produces a smooth, interpretable
    spectrum of behaviors---from aggressive all-green corridors that
    minimize EV travel time to conservative strategies that limit
    disruption to civilian traffic---giving dispatchers a single-knob
    interface to balance competing objectives.

  \item \textbf{Multi-Agent Decision Transformer (MADT) with graph
    attention for scalable intersection coordination.}  We extend the
    single-agent DT to a decentralized multi-agent architecture in
    which each intersection agent maintains its own return-conditioned
    transformer while exchanging spatial context with neighbors through
    graph attention network (GAT) layers~\cite{velickovic2018graph},
    enabling scalable coordination across networks of arbitrary
    topology.

  \item \textbf{Comprehensive empirical comparison on LightSim.}  We
    evaluate the proposed single-agent DT and MADT against online RL
    baselines (PPO, DQN), fixed-timing EVP, and state-of-the-art
    multi-agent traffic signal controllers (PressLight, CoLight) on
    grid networks of varying size and a real-world-inspired corridor,
    demonstrating competitive or superior EV travel-time reduction with
    strictly offline training.
\end{enumerate}

The remainder of this paper is organized as follows.
\Cref{sec:related} reviews related work.  \Cref{sec:formulation}
formalizes the EV corridor problem.  \Cref{sec:method} presents the
single-agent DT and MADT architectures.  \Cref{sec:setup} describes
the experimental setup, and \Cref{sec:results} reports results.
\Cref{sec:conclusion} concludes.

% ============================================================
% II. RELATED WORK
% ============================================================
\section{Related Work}
\label{sec:related}

% ------------------------------------------------------------
\subsection{Emergency Vehicle Preemption}
\label{sec:related:evp}

Traditional EV preemption systems detect an approaching EV and
immediately switch the downstream intersection to a green phase,
typically via optical (Opticom) or GPS-based priority
requests~\cite{qin2012control}.  While effective at isolated
intersections under low traffic, these reactive strategies ignore
network-level interactions: forcing green at one intersection may
propagate queues that block the EV at subsequent
intersections~\cite{wang2022ev}.

Several recent works apply \emph{online} deep RL to EV preemption.
\todo{Expand with 3--4 specific papers: cite Guo et al., Xu et al.,
and others that apply DQN/PPO/A2C to EV signal priority.  Discuss
their limitations: online training, single-intersection focus, lack of
return-conditioning.}

A common limitation across these approaches is the reliance on online
interaction---millions of simulation steps with exploratory
actions---and the absence of a mechanism for post-training behavior
modulation.

% ------------------------------------------------------------
\subsection{Decision Transformer and Offline RL}
\label{sec:related:dt}

Offline RL~\cite{levine2020offline} trains policies exclusively from a
fixed dataset of previously collected transitions, removing the need
for online exploration.  Conservative Q-Learning
(CQL)~\cite{kumar2020conservative} and Batch-Constrained deep
Q-learning (BCQ)~\cite{fujimoto2019off} regularize value estimates to
avoid overestimation on out-of-distribution actions.

The Decision Transformer~\cite{chen2021decision} recasts offline RL as
autoregressive sequence modeling: given a trajectory of
(return-to-go, state, action) tokens, a GPT-style
transformer~\cite{vaswani2017attention} predicts the next action
conditioned on a desired future return.  Trajectory
Transformer~\cite{janner2021offline} extends the idea by modeling full
state--action--reward sequences, enabling beam-search planning.

\todo{Discuss Gato, multi-game DT, and any 2023--2025 DT variants
that are relevant.  Emphasize that none have been applied to traffic
signal control or EV preemption.}

% ------------------------------------------------------------
\subsection{Multi-Agent RL for Traffic Signal Control}
\label{sec:related:marl}

Multi-agent RL (MARL) for traffic signals has progressed from
independent learners~\cite{chu2019multi} through parameter-sharing
schemes~\cite{chen2020toward} to coordination-graph
methods~\cite{wei2019colight}.  PressLight~\cite{wei2019presslight}
incorporates the max-pressure principle into the reward design, while
CoLight~\cite{wei2019colight} introduces graph attentional networks to
capture intersection-level cooperation.

\todo{Add 2--3 more recent MARL-for-traffic papers (2022--2025).
Highlight that all existing MARL traffic methods use \emph{online}
training and none incorporate return-conditioning for post-hoc
behavior adjustment.}

% ============================================================
% III. PROBLEM FORMULATION
% ============================================================
\section{Problem Formulation}
\label{sec:formulation}

% ------------------------------------------------------------
\subsection{Network and Traffic Model}
\label{sec:formulation:network}

We model the urban traffic network as a directed graph
$\mathcal{G} = (\mathcal{V}, \mathcal{E})$, where each node
$v \in \mathcal{V}$ represents a signalized intersection and each
edge $e \in \mathcal{E}$ represents a road segment connecting two
adjacent intersections.  Traffic dynamics on each segment are governed
by the Cell Transmission Model (CTM)~\cite{daganzo1994cell,
daganzo1995cell}, which discretizes road segments into cells of length
$\Delta x$ and updates vehicle counts at each time step $\Delta t$
according to:
\begin{equation}
  n_i(t+1) = n_i(t) + q_{i-1}(t) - q_i(t),
  \label{eq:ctm}
\end{equation}
where $n_i(t)$ is the number of vehicles in cell $i$ at time $t$, and
$q_i(t)$ is the flow from cell $i$ to cell $i+1$, bounded by the
sending and receiving functions:
\begin{equation}
  q_i(t) = \min\!\bigl\{
    v_f \, n_i(t),\;
    w \bigl(n_i^{\max} - n_{i+1}(t)\bigr),\;
    Q_{\max}
  \bigr\},
  \label{eq:ctm_flow}
\end{equation}
with free-flow speed $v_f$, backward wave speed $w$, maximum cell
occupancy $n_i^{\max}$, and maximum flow $Q_{\max}$.

Each intersection $v$ is controlled by a signal plan
$\phi_v(t) \in \Phi_v$ that determines which movements receive a green
indication at time $t$.

% ------------------------------------------------------------
\subsection{Emergency Vehicle Corridor Problem}
\label{sec:formulation:ev}

An \emph{emergency vehicle corridor} is a spatiotemporal signal plan
along a route $\rho = (v_1, v_2, \ldots, v_K)$ from origin $v_1$ to
destination $v_K$ such that the EV encounters green (or preemption)
phases at each intersection upon arrival.  We formalize the corridor
optimization problem as a Markov Decision Process (MDP)
$\mathcal{M} = (\mathcal{S}, \mathcal{A}, T, R, \gamma)$:

\paragraph{State space $\mathcal{S}$}
At each decision step $t$, the state $s_t \in \mathcal{S}$ comprises:
\begin{itemize}
  \item \textbf{EV state:} current edge, position along edge, speed,
    and remaining route $\rho_{t:K}$.
  \item \textbf{Traffic state:} per-cell vehicle counts
    $\{n_i(t)\}$ and per-lane queue lengths for all intersections
    within a horizon of $H$ hops from the EV's current position.
  \item \textbf{Signal state:} current phase index and elapsed phase
    time for each intersection within the $H$-hop neighborhood.
\end{itemize}
Formally,
\begin{equation}
  s_t = \bigl(\,
    \underbrace{e_t,\, p_t,\, v_t,\, \rho_{t:K}}_{\text{EV state}},\;
    \underbrace{\{n_i(t)\}_{i \in \mathcal{N}_H},\;
                \{l_j(t)\}_{j \in \mathcal{L}_H}}_{\text{traffic state}},\;
    \underbrace{\{\phi_v(t), \tau_v(t)\}_{v \in \mathcal{V}_H}}_{\text{signal state}}
  \,\bigr),
  \label{eq:state}
\end{equation}
where $\mathcal{N}_H$, $\mathcal{L}_H$, and $\mathcal{V}_H$ denote
cells, lanes, and intersections within the $H$-hop neighborhood,
respectively.

\paragraph{Action space $\mathcal{A}$}
The action at each step consists of signal-phase selections for all
controlled intersections along the corridor:
\begin{equation}
  a_t = \bigl(\phi_{v_1}(t),\, \phi_{v_2}(t),\, \ldots,\,
         \phi_{v_K}(t)\bigr) \in \mathcal{A}
         = \Phi_{v_1} \times \Phi_{v_2} \times \cdots \times \Phi_{v_K}.
  \label{eq:action}
\end{equation}
When the route is not fixed a priori, the action additionally includes
a next-edge routing decision for the EV at each intersection.

\paragraph{Transition function $T$}
The transition $T(s_{t+1} \mid s_t, a_t)$ is governed by the CTM
dynamics (\ref{eq:ctm})--(\ref{eq:ctm_flow}), the EV's kinematic
model, and background traffic demand patterns.

\paragraph{Reward function $R$}
We define a composite reward that balances EV travel time against
civilian disruption:
\begin{equation}
  r_t = R(s_t, a_t) =
    \underbrace{-\,\alpha\;\Delta d_t^{-1}}_{\text{EV progress}}
    \;+\;
    \underbrace{-\,\beta \sum_{v \in \mathcal{V}_H}
      w_v(t)}_{\text{civilian delay}}
    \;+\;
    \underbrace{\lambda\;\mathbb{1}[\text{EV arrived}]}_{\text{completion bonus}},
  \label{eq:reward}
\end{equation}
where $\Delta d_t$ is the distance traveled by the EV in step $t$,
$w_v(t)$ is the total vehicle waiting time at intersection $v$, and
$\alpha, \beta, \lambda$ are weighting coefficients.

\paragraph{Discount factor}
We use $\gamma \in [0.95, 0.99]$ throughout.

The \emph{return} of a trajectory $\tau = (s_0, a_0, r_0, \ldots,
s_T, a_T, r_T)$ is:
\begin{equation}
  R(\tau) = \sum_{t=0}^{T} \gamma^t \, r_t,
  \label{eq:return}
\end{equation}
and the \emph{return-to-go} at step $t$ is:
\begin{equation}
  \Rtg_t = \sum_{t'=t}^{T} \gamma^{t'-t}\, r_{t'}.
  \label{eq:rtg}
\end{equation}

% ------------------------------------------------------------
\subsection{Multi-Agent Formulation}
\label{sec:formulation:marl}

For large networks it is impractical to centralize control over all
intersections.  We reformulate the problem as a Decentralized Partially
Observable MDP (Dec-POMDP)
$\langle \mathcal{I}, \{\mathcal{S}^i\}, \{\mathcal{A}^i\},
T, \{R^i\}, \{\Omega^i\}, O, \gamma \rangle$,
where $\mathcal{I} = \{1, \ldots, N\}$ indexes intersection agents.
Agent~$i$ observes a local observation $o_t^i \in \Omega^i$ derived
from its $H$-hop neighborhood and selects a local phase action
$a_t^i \in \mathcal{A}^i = \Phi_{v_i}$.

The joint transition remains governed by the global CTM dynamics.
Each agent receives a \emph{shaped local reward}:
\begin{equation}
  r_t^i =
    -\,\alpha^i\;\Delta d_t^{i}\;
    -\;\beta^i \, w_{v_i}(t)\;
    +\;\lambda^i\;\mathbb{1}[\text{EV passes } v_i],
  \label{eq:local_reward}
\end{equation}
where $\Delta d_t^{i}$ captures the EV's progress on edges adjacent to
intersection~$i$.  The per-agent return-to-go $\Rtg_t^i$ is computed
analogously to~\eqref{eq:rtg}.

% ============================================================
% IV. METHODOLOGY
% ============================================================
\section{Methodology}
\label{sec:method}

% ------------------------------------------------------------
\subsection{Decision Transformer for EV Corridor}
\label{sec:method:dt}

\subsubsection{Architecture}
We adapt the Decision Transformer~\cite{chen2021decision} to the EV
corridor MDP defined in \Cref{sec:formulation:ev}.  The model
autoregressively predicts actions conditioned on a desired
return-to-go:
\begin{equation}
  a_t = \text{DT}_\theta\!\bigl(
    \Rtg_t,\, s_t,\, a_{t-1},\,
    \Rtg_{t-1},\, s_{t-1},\, \ldots,\,
    \Rtg_{t-C+1},\, s_{t-C+1}
  \bigr),
  \label{eq:dt}
\end{equation}
where $C$ is the context length (number of past time steps) and
$\theta$ denotes model parameters.

Each modality is embedded into a $d$-dimensional token:
\begin{itemize}
  \item \textbf{Return-to-go} $\Rtg_t \in \R$ is projected via a
    linear layer.
  \item \textbf{State} $s_t$ is encoded by a small MLP that
    concatenates the EV, traffic, and signal sub-vectors.
  \item \textbf{Action} $a_t$ is embedded via a learned embedding
    table (discrete phases) or linear projection (continuous).
\end{itemize}
All tokens share a learned positional encoding indexed by time step.
The token sequence is processed by a causal (GPT-style) transformer
with $L$~layers, $H_{\text{attn}}$~attention heads, and hidden
dimension~$d$.

\begin{figure}[t]
  \centering
  % \includegraphics[width=\columnwidth]{figures/dt_architecture.pdf}
  \fbox{\parbox{0.92\columnwidth}{\centering\vspace{2.5cm}
    \textbf{Figure placeholder:} Single-agent DT architecture.\\
    Input sequence of $(\Rtg, s, a)$ tokens $\to$ causal transformer
    $\to$ predicted action $a_t$.
  \vspace{0.3cm}}}
  \caption{Single-agent Decision Transformer architecture for EV
    corridor optimization.  The model receives interleaved
    return-to-go, state, and action tokens over a context window of
    $C$ time steps and predicts the next action.}
  \label{fig:dt_arch}
\end{figure}

\subsubsection{Return-Conditioning Mechanism}
At training time the return-to-go tokens $\Rtg_t$ are computed from
the dataset trajectory via~\eqref{eq:rtg}.  At inference time, the
dispatcher specifies a \emph{target return} $G^\star$ that seeds the
initial token $\Rtg_0 = G^\star$; subsequent return-to-go values are
updated online as rewards are observed:
\begin{equation}
  \Rtg_{t+1} = \Rtg_t - r_t.
  \label{eq:rtg_update}
\end{equation}
A high $G^\star$ biases the model toward aggressive green-wave
corridors (minimizing EV delay at the cost of civilian traffic), while
a lower $G^\star$ yields conservative strategies with smaller
disruption.

\subsubsection{Training Procedure}
The model is trained offline on a dataset
$\mathcal{D} = \{\tau^{(j)}\}_{j=1}^{M}$ of $M$ trajectories
collected by a mixture of policies (random, heuristic, and
partially-trained RL agents) to ensure broad coverage of the
return distribution.  Training minimizes the cross-entropy loss for
discrete actions (or MSE for continuous):
\begin{equation}
  \mathcal{L}(\theta) = -\;\E_{(\Rtg,\, s,\, a) \sim \mathcal{D}}
    \bigl[\log \pi_\theta(a_t \mid \Rtg_{t:t-C+1},\, s_{t:t-C+1},\,
    a_{t-1:t-C+1})\bigr].
  \label{eq:dt_loss}
\end{equation}

% ------------------------------------------------------------
\subsection{Multi-Agent Decision Transformer}
\label{sec:method:madt}

\subsubsection{Architecture Overview}
In the MADT~\cite{su2026madt}, each intersection agent~$i$ maintains a
local DT with shared parameters $\theta$.  Before the causal
transformer layers, agent embeddings are enriched with
\emph{spatial context} from neighboring agents via a Graph Attention
Network (GAT)~\cite{velickovic2018graph} operating on the intersection
graph $\mathcal{G}$.

Let $\bm{h}_t^i \in \R^d$ denote agent $i$'s embedding at time $t$
after the token-embedding layer.  The GAT computes:
\begin{align}
  \alpha_{ij}^t &= \frac{
    \exp\!\bigl(\text{LeakyReLU}(\bm{a}^\top
    [\bm{W}\bm{h}_t^i \,\|\, \bm{W}\bm{h}_t^j])\bigr)
  }{
    \sum_{k \in \mathcal{N}(i)}
    \exp\!\bigl(\text{LeakyReLU}(\bm{a}^\top
    [\bm{W}\bm{h}_t^i \,\|\, \bm{W}\bm{h}_t^k])\bigr)
  }, \label{eq:gat_attn} \\[4pt]
  \tilde{\bm{h}}_t^i &= \sigma\!\Bigl(
    \sum_{j \in \mathcal{N}(i)} \alpha_{ij}^t\;\bm{W}\bm{h}_t^j
  \Bigr), \label{eq:gat_agg}
\end{align}
where $\bm{W} \in \R^{d \times d}$ and $\bm{a} \in \R^{2d}$ are
learnable parameters, $\|$ denotes concatenation, and
$\mathcal{N}(i)$ is the set of neighbors of intersection~$i$ in
$\mathcal{G}$.  We use multi-head attention with $K$ heads and
concatenate their outputs.

The enriched embedding $\tilde{\bm{h}}_t^i$ replaces the original
state token and is fed into the shared causal transformer, yielding
agent $i$'s action prediction.

\begin{figure}[t]
  \centering
  % \includegraphics[width=\columnwidth]{figures/madt_architecture.pdf}
  \fbox{\parbox{0.92\columnwidth}{\centering\vspace{3.0cm}
    \textbf{Figure placeholder:} MADT architecture.\\
    Per-agent token embeddings $\to$ GAT spatial aggregation $\to$
    shared causal transformer $\to$ per-agent action predictions.
  \vspace{0.3cm}}}
  \caption{Multi-Agent Decision Transformer (MADT) architecture.
    Each intersection agent embeds its local $(\Rtg^i, o^i, a^i)$
    tokens, exchanges spatial context with neighbors through GAT
    layers, and then feeds the enriched embeddings into a shared
    causal transformer for action prediction.}
  \label{fig:madt_arch}
\end{figure}

\subsubsection{Per-Agent Return Conditioning}
Each agent~$i$ is conditioned on its own local return-to-go
$\Rtg_t^i$, derived from the local reward~\eqref{eq:local_reward}.
At inference, the dispatcher may broadcast a uniform target
$G^\star$ to all agents or assign differentiated targets depending on
proximity to the EV's predicted path.

\subsubsection{Training}
MADT is trained with the same offline loss~\eqref{eq:dt_loss},
applied independently per agent but with shared parameters, on a
multi-agent dataset $\mathcal{D} = \{(\tau^{(j),1}, \ldots,
\tau^{(j),N})\}_{j=1}^M$.  The GAT parameters are trained
end-to-end.

% ------------------------------------------------------------
\subsection{Offline Data Collection Strategy}
\label{sec:method:data}

High-quality offline data with broad return coverage is essential for
return-conditioned models.  We collect $\mathcal{D}$ in three phases:

\begin{enumerate}
  \item \textbf{Random policy} ($\sim$30\% of dataset): uniform random
    phase selections, yielding low-return trajectories.
  \item \textbf{Heuristic policy} ($\sim$40\%): a fixed-timing EVP
    heuristic that grants green to the EV's current approach whenever
    the EV is within $D$ meters, producing medium-return trajectories.
  \item \textbf{Trained RL policy} ($\sim$30\%): a partially-converged
    PPO agent~\cite{schulman2017proximal} providing high-return
    trajectories.
\end{enumerate}

\todo{Add details on dataset sizes, trajectory lengths, and any data
augmentation (e.g., hindsight return relabeling).}

% ============================================================
% V. EXPERIMENTAL SETUP
% ============================================================
\section{Experimental Setup}
\label{sec:setup}

% ------------------------------------------------------------
\subsection{Simulation Environment}
\label{sec:setup:sim}

All experiments are conducted in \textbf{LightSim}~\cite{lightsim2024},
a lightweight, GPU-accelerated traffic signal simulator designed for
RL research.  LightSim implements CTM-based traffic dynamics and
supports configurable intersection topologies, demand patterns, and
EV injection.

\todo{Add LightSim version, key simulation parameters ($\Delta t$,
$\Delta x$, demand levels), and any modifications made for EV
simulation.}

% ------------------------------------------------------------
\subsection{Network Scenarios}
\label{sec:setup:networks}

We evaluate on three network configurations:
\begin{enumerate}
  \item \textbf{$4 \times 4$ grid:} 16 signalized intersections with
    uniform demand; EV travels from southwest corner to northeast
    corner.
  \item \textbf{$6 \times 6$ grid:} 36 intersections, same demand
    structure, to test scalability.
  \item \textbf{Real-world corridor:} \todo{Describe real-world-inspired
    network (e.g., based on a segment of city X with Y intersections).}
\end{enumerate}

% ------------------------------------------------------------
\subsection{Baselines}
\label{sec:setup:baselines}

\begin{itemize}
  \item \textbf{Fixed-Time EVP (FT-EVP):} standard preemption that
    grants green to the EV approach when the EV is within a detection
    range of $D = 300\,\text{m}$.
  \item \textbf{DQN}~\cite{mnih2015human}: online deep Q-network
    trained with the same reward~\eqref{eq:reward}.
  \item \textbf{PPO}~\cite{schulman2017proximal}: online proximal
    policy optimization.
  \item \textbf{PressLight}~\cite{wei2019presslight}: max-pressure RL
    for multi-agent signal control (no EV-specific reward; green given
    to EV lane by max-pressure criterion).
  \item \textbf{CoLight}~\cite{wei2019colight}: graph-attention MARL
    baseline.
\end{itemize}

\todo{Confirm final baseline list; add CQL~\cite{kumar2020conservative}
as an offline RL baseline for a fairer comparison.}

% ------------------------------------------------------------
\subsection{Evaluation Metrics}
\label{sec:setup:metrics}

\begin{itemize}
  \item \textbf{EV Travel Time (ETT):} total time for the EV to
    traverse from origin to destination.
  \item \textbf{Average Civilian Delay (ACD):} mean increase in
    travel time for non-EV vehicles relative to free-flow conditions.
  \item \textbf{Throughput:} total vehicles served during the episode.
  \item \textbf{EV Stops:} number of times the EV speed drops below
    $1\,\text{m/s}$.
\end{itemize}

% ------------------------------------------------------------
\subsection{Implementation Details}
\label{sec:setup:impl}

\todo{Fill in: transformer hyperparameters ($L$, $H_{\text{attn}}$,
$d$, $C$), GAT configuration ($K$ heads, layers), learning rate,
batch size, number of training steps, hardware (GPU model), training
wall-clock time.  Reference the codebase.}

% ============================================================
% VI. RESULTS AND DISCUSSION
% ============================================================
\section{Results and Discussion}
\label{sec:results}

% ------------------------------------------------------------
\subsection{Main Results}
\label{sec:results:main}

\begin{table*}[t]
  \centering
  \caption{Main results on the $4\times4$ grid network.  Best results
    in \textbf{bold}, second best \underline{underlined}.
    \todo{Fill with experimental data.}}
  \label{tab:main}
  \begin{tabular}{lcccc}
    \toprule
    \textbf{Method} & \textbf{ETT (s) $\downarrow$}
      & \textbf{ACD (s) $\downarrow$}
      & \textbf{Throughput $\uparrow$}
      & \textbf{EV Stops $\downarrow$} \\
    \midrule
    FT-EVP          & \todo{} & \todo{} & \todo{} & \todo{} \\
    DQN             & \todo{} & \todo{} & \todo{} & \todo{} \\
    PPO             & \todo{} & \todo{} & \todo{} & \todo{} \\
    PressLight      & \todo{} & \todo{} & \todo{} & \todo{} \\
    CoLight         & \todo{} & \todo{} & \todo{} & \todo{} \\
    \midrule
    DT (ours)       & \todo{} & \todo{} & \todo{} & \todo{} \\
    MADT (ours)     & \todo{} & \todo{} & \todo{} & \todo{} \\
    \bottomrule
  \end{tabular}
\end{table*}

\todo{Discuss main results.  Highlight (1) DT matches or beats online
RL despite offline-only training, (2) MADT further improves through
spatial coordination, (3) both methods produce far fewer EV stops than
FT-EVP.}

% ------------------------------------------------------------
\subsection{Return-Conditioning Analysis}
\label{sec:results:return}

\begin{figure}[t]
  \centering
  % \includegraphics[width=\columnwidth]{figures/return_sweep.pdf}
  \fbox{\parbox{0.92\columnwidth}{\centering\vspace{3.0cm}
    \textbf{Figure placeholder:} Return-conditioning sweep.\\
    $x$-axis: target return $G^\star$; $y$-axis: achieved ETT and ACD.
    Curves show smooth trade-off.
  \vspace{0.3cm}}}
  \caption{Effect of target return $G^\star$ on EV travel time (ETT)
    and average civilian delay (ACD).  Higher targets produce faster
    EV corridors at the cost of increased civilian delay, providing a
    single-knob dispatch interface.}
  \label{fig:return_sweep}
\end{figure}

\todo{Present return-conditioning sweep results.  Show that varying
$G^\star$ produces a smooth Pareto front between ETT and ACD.  Compare
against PPO (which has no such knob and would require retraining with
different reward weights).}

% ------------------------------------------------------------
\subsection{Scalability Analysis}
\label{sec:results:scale}

\begin{figure}[t]
  \centering
  % \includegraphics[width=\columnwidth]{figures/scalability.pdf}
  \fbox{\parbox{0.92\columnwidth}{\centering\vspace{2.5cm}
    \textbf{Figure placeholder:} Scalability analysis.\\
    $x$-axis: grid size ($4\times4$, $6\times6$, \ldots);
    $y$-axis: ETT.  Lines for each method.
  \vspace{0.3cm}}}
  \caption{Scalability of different methods as network size increases.
    MADT maintains competitive performance on larger grids due to
    localized graph-attention coordination.}
  \label{fig:scalability}
\end{figure}

\todo{Compare performance degradation across grid sizes.  Show that
MADT with GAT scales more gracefully than centralized DT and DQN.}

% ------------------------------------------------------------
\subsection{Ablation Studies}
\label{sec:results:ablation}

\begin{table}[t]
  \centering
  \caption{Ablation study on the $4\times4$ grid.
    \todo{Fill with experimental data.}}
  \label{tab:ablation}
  \begin{tabular}{lcc}
    \toprule
    \textbf{Variant} & \textbf{ETT (s)} & \textbf{ACD (s)} \\
    \midrule
    MADT (full)              & \todo{} & \todo{} \\
    \quad w/o GAT            & \todo{} & \todo{} \\
    \quad w/o return cond.   & \todo{} & \todo{} \\
    \quad context $C = 5$    & \todo{} & \todo{} \\
    \quad context $C = 10$   & \todo{} & \todo{} \\
    \quad context $C = 20$   & \todo{} & \todo{} \\
    \quad low-quality data   & \todo{} & \todo{} \\
    \bottomrule
  \end{tabular}
\end{table}

\todo{Discuss ablations:  (1) removing GAT hurts multi-agent
coordination; (2) removing return conditioning reduces controllability;
(3) context length has diminishing returns beyond $C=10$;
(4) dataset quality significantly impacts DT performance---mixed
datasets outperform expert-only and random-only.}

\begin{figure}[t]
  \centering
  % \includegraphics[width=\columnwidth]{figures/learning_curves.pdf}
  \fbox{\parbox{0.92\columnwidth}{\centering\vspace{2.5cm}
    \textbf{Figure placeholder:} Training loss / evaluation curves.\\
    Offline DT loss convergence + online baseline reward curves.
  \vspace{0.3cm}}}
  \caption{Training convergence.  Left: DT/MADT offline training loss.
    Right: Online RL (PPO, DQN) reward curves for reference.}
  \label{fig:learning}
\end{figure}

\begin{figure}[t]
  \centering
  % \includegraphics[width=\columnwidth]{figures/corridor_vis.pdf}
  \fbox{\parbox{0.92\columnwidth}{\centering\vspace{3.0cm}
    \textbf{Figure placeholder:} Network visualization.\\
    Heatmap of signal phases and EV trajectory on the grid network
    under different $G^\star$ settings.
  \vspace{0.3cm}}}
  \caption{Visualization of EV corridors produced by MADT under
    different target returns.  Green cells indicate preempted phases;
    the EV path is shown in red.}
  \label{fig:corridor_vis}
\end{figure}

% ============================================================
% VII. CONCLUSION
% ============================================================
\section{Conclusion}
\label{sec:conclusion}

This paper introduced a return-conditioned framework for emergency
vehicle corridor optimization based on the Decision Transformer.
By framing the corridor problem as offline sequence modeling, the
proposed approach eliminates the need for exploratory online
interaction---a critical advantage in safety-sensitive EV
operations---while providing dispatchers with an intuitive
target-return knob to adjust urgency in real time.  The Multi-Agent
Decision Transformer (MADT) extension incorporates graph attention
layers to enable decentralized, scalable coordination across
intersections of arbitrary network topology.

\todo{Summarize key quantitative findings once experiments are
complete.}

Limitations and future directions include:
\begin{itemize}
  \item \todo{Discuss sim-to-real transfer challenges.}
  \item \todo{Discuss extension to multiple simultaneous EVs.}
  \item \todo{Discuss integration with real-time traffic prediction
    and V2X communication.}
\end{itemize}

% ============================================================
% REFERENCES
% ============================================================
\bibliographystyle{IEEEtran}
\bibliography{references}

\end{document}
